\section{Related Work}
\label{sec:related_work}

\textbf{Code Chunking for RAG:} Line-based chunking can disrupt semantic boundaries~\cite{zhang2025cast}. \castname~\cite{zhang2025cast} addresses this by recursively partitioning AST nodes and merging siblings under size limits, reporting gains on code-completion and issue-resolution benchmarks. However, subsequent controlled experiments in code completion show that structure-aware chunking does not reliably outperform sliding windows~\cite{wu2026chunking}, and line-based slicing matches syntax-aware splitting across compute budgets~\cite{galimzyanov2025practical}. While those studies target code completion, we evaluate human-labeled requirement-to-code evidence retrieval and orthogonally control context headers.

\textbf{Dense Retrieval and Reranking:} Neural code search embeds natural-language queries and source code into shared representations~\cite{feng2020codebert,guo2022unixcoder} (benchmarked on CodeSearchNet~\cite{husain2019codesearchnet}). Modern dense models like BGE-M3~\cite{chen2024bgem3} enable multilingual semantic matching; two-stage pipelines integrate cross-encoder rerankers that jointly attend across query and chunk tokens~\cite{hu2023revisiting,costv2024apsec} to overcome single-vector limits. We adopt this architecture to examine interactions between chunk geometry, headers, and cross-encoders.

\textbf{Context Enrichment:} Isolated code chunks frequently lack enclosing class and package context. Prior approaches enrich completion prompts with repository context iteratively (RepoCoder~\cite{zhang2023repocoder}) or benchmark cross-file dependencies (RepoBench~\cite{liu2024repobench}, CrossCodeEval~\cite{ding2023crosscodeeval}). We investigate a lightweight alternative: prepending standardized context headers to individual chunks prior to retrieval.

\textbf{Requirements Traceability and Skill Mining:} Tracing natural-language requirements to code has deep roots in software engineering~\cite{antoniol2002traceability,gao2024triad}. Separately, repository mining classifies developer expertise from framework usage~\cite{montandon2019identifying}. We adopt the traceability perspective: treating required technical skills as requirements and retrieving concrete code evidence, without assessing individual developers.
